\documentclass{article}
\usepackage[a4paper,margin=10mm,onecolumn,bottom=2cm,footskip=1cm]{geometry}
\usepackage{lmodern}
\usepackage{minted}
\usepackage[T1]{fontenc}
\usepackage{tfrupee}
\setlength{\headheight}{1cm}%set height of header box
\usepackage{graphicx}% Required for inserting images
\setlength{\headsep}{4mm}
\usepackage{setspace}
\setstretch{1.0}%Space Between lines
\usepackage{csquotes}
\usepackage{cite}
\usepackage{amsmath,amssymb,amsfonts,amsthm}
\usepackage{enumitem}
\usepackage{adjustbox}
\usepackage{fancyhdr}
\usepackage{lastpage}
\usepackage{romannum}
\usepackage{listings}
\input{listings-python.prf}  % defines the python-idle-code style
\usepackage{textcomp} % Needed for straight quotes 
\usepackage{minted}
\usemintedstyle[c]{colorful}

\lstdefinestyle{pycode}{
 basicstyle=\normalsize\ttfamily, % size of the fonts for the code               
 language={Python},style={python-idle-code},
 showstringspaces=false,
 tabsize=4,
 upquote=true,   % Requires textcomp
 linewidth=7in  % This just fits 80 characters at 10pt 
}

\definecolor{mygreen}{rgb}{0,0.6,0}
\definecolor{mygray}{rgb}{0.5,0.5,0.5}
\definecolor{mymauve}{rgb}{0.58,0,0.82}
\definecolor{terminalbgcolor}{HTML}{330033}
\definecolor{terminalrulecolor}{HTML}{000099}
\lstdefinestyle{terminal}{
	backgroundcolor=\color{terminalbgcolor},
	basicstyle=\color{white}\fontfamily{fvm}\footnotesize\selectfont,
	breakatwhitespace=false,  
	breaklines=true,
	captionpos=b,
	commentstyle=\color{mygreen},
	deletekeywords={...},
	escapeinside={\%*}{*)},
	extendedchars=true,
	frame=single,
	keepspaces=true,
	keywordstyle=\color{blue},
	%language=none,
	morekeywords={*,...},
	numbers=none,
	numbersep=5pt,
    %framerule=2pt,
	numberstyle=\color{mygray}\tiny\selectfont,
	rulecolor=\color{terminalrulecolor},
	showspaces=false,
	showstringspaces=false,
	showtabs=false,
	stepnumber=2,
	stringstyle=\color{mymauve},
	tabsize=2
}

\definecolor{myred}{rgb}{0.92,0.25,0.20}

%\title{\HUGE GATE 2025 \\ \vspace{0.5cm} AGRICULTURAL ENGINEERING (AG)}
%\author{Chamarthi Krishna Madhur}
%\date{\today}

\begin{document}
\begin{titlepage}
    \begin{center}
        \vfill
        \vspace*{5cm}
        \Huge\textbf{GATE 2025 \\ \vspace{0.5cm} AGRICULTURAL ENGINEERING (AG)}\\
        \vspace{2cm}
        \huge \textbf{{Answers to selected questions}}\\
        \vspace{4cm}
        \huge Chamarthi Krishna Madhur\\
        \vspace{1cm}
        
        \today
        
    \end{center}
\end{titlepage}

\pagenumbering{arabic}
\pagestyle{fancy}

\newpage
\fancyhf{} % clear existing header/footer entries
\fancyhead[R]{{\textbf{Agricultural Engineering (AG)}}}
%\renewcommand{\headrule}{\color{magenta}\hrule height 1.5pt}
\fancyfoot[C]{Page \textbf{\thepage} out of \textbf{\pageref{LastPage}} }
\fancyfoot[L]{Organizing Institute: IIT Roorkee}
\renewcommand{\footrulewidth}{0.4pt}
%\renewcommand{\footrule}{\color{magenta}\hrule height 1.5pt}





%\section*{\underline {General Aptitude}}
 \subsection*{Q.6 – Q.10 Carry TWO marks Each}
 
 \hfill\begin{minipage}{\linewidth}
 \begin{enumerate}
    \item[Q.8] A fair six-faced dice, with the faces labelled ‘1’, ‘2’, ‘3’, ‘4’, ‘5’, and ‘6’, is rolled
    thrice. What is the probability of rolling ‘6’ exactly once?
    \item[(A)] $\dfrac{75}{216}$ 
    \item[(B)] $\dfrac{1}{6}$
    \item[(C)] $\dfrac{1}{18}$
    \item[(D)] $\dfrac{25}{216}$
    %\item[SOL] NEED TO WRITE SOLUTION!!!!
\end{enumerate}
\end{minipage}\bigskip\\
\underline{Solution}\\\\
     Step 1: Finding the probability mass function(PMF) of a single dice roll:\\\\
     The sample space is S $\in \{1,2,3,4,5,6\}$ \\
     Let us define a variable $Y_1$ as follows: 
     \begin{align*}
         Y_1=
            \begin{cases}
                1&X_1=6\\
                0&X_1\neq6
            \end{cases}
     \end{align*}
     Now, The PMF is defined as:\bigskip\\
        \begin{align*}
            P_{Y_1}(k)=
                \begin{cases}
                    \frac{1}{6} &  k=1 \\
                    \frac{5}{6} & k=0\\
                \end{cases}\\\\
            \text{Let p,q be $\frac{1}{6}$ and $\frac{5}{6}$ respectively.}
        \end{align*}\\
        
    \noindent Step 2: Applying Z-Transform to find the required PMF for Y \\\\
    Let us define Y\textsubscript{2} and X\textsubscript{3} as the random variables for the 2\textsuperscript{nd} and 3\textsuperscript{rd} dice roll, respectively. Let Y denote the random variable for the combined triple dice roll.\medskip\\
    Now, we can write Y=Y\textsubscript{1}+Y\textsubscript{2}+Y\textsubscript{3}.\medskip\\
    Writing the Z-Transform for Y\textsubscript{1}\\
        \begin{align*}
        M_{Y_1}(Z)=E[Z^{-Y_1}]&=\sum\, P_{Y_1}(k)\,Z^{-k}\tag{1}\\
                           &=(P_{Y_1}(1)Z^{-1} +P_{Y_1}(0)Z^{-0})\\
                           &=pZ^{-1}+q\tag{2}
        \end{align*}
    We know that $Y_1,Y_2,Y_3$ are independent events. So, we can write $M_Y(Z)=M_{Y_1}(Z)\,\, M_{Y_2}(Z)\,\,M_{Y_3}(Z)$.\\
    Therefore, 
    \begin{align*}
        M_Y(Z)&=(pZ^{-1}+q)^3\tag{From 2}.\\
              &=\sum_{k=0}^3{^3}C_k\,\,q^{3-k}\,\, (pZ^{-1})^k\\    
              &=\sum P_Y(k)Z^{-k}\tag{From 1}\\
        \text{So, we can say:}\\
            P_Y(k)&=^3C_k\,\,q^{3-k}\,\,p^k\tag{3}
    \end{align*}
    The above situation is called as a binomial distribution.\\
    \noindent Step 3: Using binomial distribution to get required probability
    We need the situation where p appears exactly once. So, we consider only the term where the power of p is 1.\\
    The required term is :  $^3C_1\,\,p\,\,q^2$.\\
    Substituting the values of p and q, we get:\\
    \begin{center}
        $^3C_1\,\frac{1}{6}\,\,\left(\frac{5}{6}\right)^2$\medskip\\
        $=\dfrac{75}{216}$
    \end{center}
    Therefore, the required answer is $\dfrac{75}{216}$\medskip\\
\large\underline{\textbf{C Implementation}}\medskip\\
\begin{minted}{c}
#include <stdio.h>
#include<stdlib.h>
#include <time.h>
void roll(int*b);
int main(void){
  int hit=0;
  int fail=0;
  srand(time(NULL));
  int a[]={0,0,0};
  for (int i=0;i<1000000000;i++){
  roll(a);
  if ((a[0]==6 && a[1]!=6 && a[2]!=6)||(a[0]!=6 && a[1]!=6 && a[2]==6)||(a[0]!=6 && a[1]==6 && a[2]!=6)){
  //printf("hit");
  hit+=1;
  //printf("%d %d %d",a[0],a[1],a[2]);
  }
  else{
  fail+=1;
  }
  }
  printf("hits %d, Fail %d\n",hit,fail);
  printf("\nratio=%lf\n",(double)(hit)/(hit+fail)); 
  
  return 0;
}
void roll(int*b){
b[0]=rand()%6+1;
b[1]=rand()%6+1;
b[2]=rand()%6+1;
}       
\end{minted}
\vspace{1cm}
The terminal output is shown below:\medskip\\
\begin{lstlisting}[style=terminal]
madhur@madhur-Dell-15-DC15250:~$ gcc test.c
madhur@madhur-Dell-15-DC15250:~$ ./a.out
hits 347221618, Fail 652778382

ratio=0.347222
\end{lstlisting}
The answer$\left(\dfrac{75}{216}\right)$ is 0.34722222.. so, we can say that our answer is correct.\\
 \subsection*{Q.11 – Q.35 Carry ONE mark Each}
 
    \hfill\begin{minipage}{\linewidth}
\begin{enumerate}
    \item[Q.11] If a, b, and c are the roots of $2x^3 - 3x^2 + px - 1=0$ and sum of the two roots is 1, the value of $p$ is: 
    \item[(A)] 2
    \item[(B)] 1
    \item[(C)] 3
    \item[(D)] $\frac{1}{3}$
    %\item[SOL]NEED TO WRITE SOL!!!\\
\end{enumerate}
    \underline{Solution}\medskip\\
    Let the roots whose sum is 1 be a and b.\\
    Then, by using the properties of coefficients of an equation, we write:\medskip\\
    \begin{align*}
        -\left(\frac{-3}{2}\right)&=a+b+c\\
                               \frac{3}{2} &=1+c\\
                               c&=\frac{1}{2} \tag{1}
    \end{align*}
    And
    \begin{align*}
        abc&=-\left(\frac{-1}{2}\right)\\
        ab\,\frac{1}{2}&=\frac{1}{2}\tag{From (1)}\\ 
        ab&=1\tag{2}
    \end{align*}
    Also,
    \begin{align*}
        ab+bc+ac&=\left(\frac{p}{2}\right)\\
        1+c(a+b)&=\frac{p}{2}\tag{From (2)}\\
        1+\frac{1}{2}(1)&=\frac{p}{2}\tag{From (1)}\\
        3=p
    \end{align*}
    Therefore, the value of p is 3.
    
    
    \end{minipage}\bigskip

\hfill\begin{minipage}{\linewidth}
    \begin{enumerate}
    \item[Q.12] If the matrix $\begin{bmatrix} $x$&$-3$\\ $2$&$x-5$ \end{bmatrix}$ is singular, sum of the values of $x$ is:
    \item[(A)] 5
    \item [(B)]6
    \item [(C)] 0
    \item [(D)]$-1$
    %\item[SOL]NEED TO WRITE SOL!!!
\end{enumerate}
\end{minipage}\bigskip\\
\underline{Solution}\medskip\\
A matrix is said to be singular if its row echelon form has a row full of zeroes.\medskip\\
Step 1: Finding row echelon form of the given matrix:\medskip\\

\begin{center}
      $\begin{bmatrix}
        x&-3\\
        2&x-5 
        \end{bmatrix}$\bigskip\medskip\\
        $\begin{bmatrix}
        x&-3\\
        1&\frac{(x-5)}{2} 
        \end{bmatrix}$\bigskip\medskip\\
        $\begin{bmatrix}
        x-[(x)1]&-3-[(x)\frac{(x-5)}{2}]\\
        1&\frac{(x-5)}{2} 
        \end{bmatrix}$\bigskip\medskip\\
        $\begin{bmatrix}
        0&\frac{(-6-x^2+5x)}{2}\\
        1&\frac{(x-5)}{2} 
        \end{bmatrix}$\bigskip\medskip\\
        $\begin{bmatrix}
        1&\frac{(x-5)}{2}\\
        0&\frac{(-6-x^2+5x)}{2}
        \end{bmatrix}$\\  
\end{center}
\noindent Step 2: Making the bottom row to be full of zeroes:\medskip\\
We need to equate $\dfrac{(-6-x^2+5x)}{2}$ to zero.\\
We get 
\begin{align*}
    -6-x^2+5x=0\\
    x^2-5x+6=0
\end{align*}
So, the sum of values of x is 5.\medskip\\
Since it is not possible to display 'x' and 'x-5' explicitly in the code, a numerical approach was taken:\\\\
The following code converts a numerical matrix into its row echelon form.\\\\
We can verify that the values of ($x=2,3$) give us a matrix with a row of zeroes\\\\
\large \underline{\textbf{C Implementation}}\\
\begin{minted}[frame=single, framesep=10pt]{c}
#include <stdio.h>
#include <math.h>
#define ROWS 2
#define COLS 2

void printmtrx(double mtrx[ROWS][COLS]);
void convertToEchelon(double mtrx[ROWS][COLS]);

int main() {
    
    double mtrxa[ROWS][COLS] = {//for x=2
        { 2, 2},
        {-3,-3}
        
    };
    double mtrxb[ROWS][COLS]={//for x=3
      {3,-3},
      {2,-2}
      
    };
    printf("Original matrix A :\n");
    printmtrx(mtrxa);
    printf("Original matrix B :\n");
    printmtrx(mtrxb);
    convertToEchelon(mtrxa);
    convertToEchelon(mtrxb);
    

    printf("Echelon Form a:\n");
    printmtrx(mtrxa);
    printf("\n Echelon Form b : \n");
    printmtrx(mtrxb);
    return 0;
}
void printmtrx(double mtrx[ROWS][COLS]) {
    for (int i = 0; i < ROWS; i++) {
        for (int j = 0; j < COLS; j++) {
            printf("%0.2f ", mtrx[i][j]);
        }
        printf("\n");
    }
    printf("\n");
}
void convertToEchelon(double  mtrx[ROWS][COLS]) {
    int lead = 0; // Tracks the current pivot column
    for (int r = 0; r < ROWS; r++) {
        if (lead >= COLS) {
            return;
        }
        // Finding non zero element in column 
        int i = r;
        while (mtrx[i][lead] == 0) {
            i++;
            if (i == ROWS) {
                i = r;
                lead++;
                if (lead == COLS) {
                    return;
                }
            }
        }
        // Swap current row 'r' with pivot row 'i'
        if (i != r) {
            for (int j = 0; j < COLS; j++) {
                double temp = mtrx[r][j];
                mtrx[r][j] = mtrx[i][j];
                mtrx[i][j] = temp;
            }
        }
        // 2. Eliminate entries below the current pivot
        for (int k = r + 1; k < ROWS; k++) {
            if (fabs(mtrx[k][lead]) > 1e-9) {
                double factor = mtrx[k][lead] / mtrx[r][lead];
                for (int j = lead; j < COLS; j++) {
                    mtrx[k][j] -= factor * mtrx[r][j];
                }
            }
        }
        
        lead++;
    }
}
\end{minted}
The terminal output is :\\
\begin{lstlisting}[style=terminal]
madhur@madhur-Dell-15-DC15250:~/LATEX_CODES$ gcc echelon.c -lm
madhur@madhur-Dell-15-DC15250:~/LATEX_CODES$ ./a.out
Original matrix A :
2.00 2.00 
-3.00 -3.00 

Original matrix B :
3.00 -3.00 
2.00 -2.00 

Echelon Form a:
2.00 2.00 
0.00 0.00 


 Echelon Form b : 
3.00 -3.00 
0.00 0.00 
\end{lstlisting}
Since both gave a row echelon form with a row of zeroes, we can conclude that $x=2,3$ are the 2 values satisfying the singular condition. So, their sum is 5.\medskip\\


\hfill\begin{minipage}{\linewidth}
    \begin{enumerate}
    \item[Q.13] $\dfrac{s}{s^2+a^2}$ is the Laplace transform of: 
    \item[(A)] cos$(at)$
    \item[(B)] sin$(at)$
    \item[(C)] sinh$(at)$
    \item[(D)] cosh$(at)$
    %\item[SOL]NEED TO WRITE SOL!!!            
    \end{enumerate}
    \end{minipage}\bigskip\\
\underline{Solution}\medskip\\
We do not know the laplace transforms of the options, so we hit-and-trial them to get the expression:\medskip\\
Option (A):\\
\begin{align*}
    \mathcal{L}\{\cos(at)\}&=\int_{0}^{\infty}\cos(at)\,e^{-st}\,dt=I\\
                          I&=\int_{0}^{\infty}\frac{d(\sin(at))}{a}\,e^{-st}\\
                           &=\frac{\sin(at)}{a}\,e^{-st}\Bigg|_0^\infty-\frac{1}{a}\int_{0}^{\infty}\sin(at)e^{-st}(-s)dt\\
                           &=0+\frac{s}{a}\int_0^\infty \frac{d(-\cos(at))}{a}\,e^{-st}\\
                           &=-(\frac{s}{a^2})\left\{\cos(at)\,e^{-st}\,\,\Bigg|_0^\infty-\int_0^\infty\cos(at)\,e^{-st}\,(-s)\right\}\\
                           &=-(\frac{s}{a^2})\left\{(0-1)-(-s)I\right\}\\
                           &=-(\frac{s}{a^2})\left\{-1+sI\right\}\\
                          I&=\frac{s}{a^2}-\frac{s^2}{a^2}(I)\\
       \left(\frac{a^2+s^2}{a^2}\right)I&=\frac{s}{a^2}\\
                          I&=\frac{s}{s^2+a^2}
\end{align*}\bigskip\\
\begin{center}
    (continued in the next page)
\end{center}
Option (B):
\begin{align*}
    \mathcal{L}\{\sin(at)\}&=\int_{0}^{\infty}\sin(at)\,e^{-st}\,dt=I\\
                          I&=\int_{0}^{\infty}\frac{d(-\cos(at))}{a}\,e^{-st}\\
                           &=\frac{-\cos(at)}{a}\,e^{-st}\Bigg|_0^\infty-\frac{1}{a}\int_{0}^{\infty}-\cos(at)e^{-st}(-s)dt\\
                           &=\frac{1}{a}-\frac{s}{a}\int_0^\infty \frac{d(\sin(at))}{a}\,e^{-st}\\
                           &=\frac{1}{a}-(\frac{s}{a^2})\left\{\sin(at)\,e^{-st}\,\,\Bigg|_0^\infty-\int_0^\infty\sin(at)\,e^{-st}\,(-s)\right\}\\
                           &=-(\frac{s}{a^2})\left\{(0-0)-(-s)I\right\}\\
                           &=-(\frac{s}{a^2})\left\{0+sI\right\}\\
                          I&=\frac{a}{a^2}-\frac{s^2}{a^2}(I)\\
       \left(\frac{a^2+s^2}{a^2}\right)I&=\frac{a}{a^2}\\
                          I&=\frac{a}{s^2+a^2}
\end{align*}\bigskip\\\\
Option (C):\\
\begin{align*}
    \mathcal{L}\{\sinh(at)\}&=\int_0^\infty \frac{e^{at}-e^{-at}}{2}\,e^{-st} dt\\
                            &=\frac{1}{2}\int_0^\infty e^{(a-s)(t)} dt -\frac{1}{2}\int_0^\infty e^{(-a-s)(t)}dt\\
                            &=\frac{1}{2}\left\{\frac{e^{(a-s)(t)}}{a-s}\right\}\Bigg|_0^\infty-\frac{1}{2}\left\{\frac{e^{(-a-s)(t)}}{-a-s}\right\}\Bigg|_0^\infty\tag{assume s$>$a for integral to be finite}\\
                            &=\frac{1}{2}\left(\frac{1}{s-a}-\frac{1}{s+a}\right)\\
                            &=\frac{1}{2}\left(\frac{2a}{s^2-a^2}\right)\\
                           I&=\frac{a}{s^2-a^2}
\end{align*}
Option (D):\\
\begin{align*}
    \mathcal{L}\{\cosh(at)\}&=\int_0^\infty \frac{e^{at}+e^{-at}}{2}\,e^{-st} dt\\
                            &=\frac{1}{2}\int_0^\infty e^{(a-s)(t)} dt +\frac{1}{2}\int_0^\infty e^{(-a-s)(t)}dt\\
                            &=\frac{1}{2}\left\{\frac{e^{(a-s)(t)}}{a-s}\right\}\Bigg|_0^\infty+\frac{1}{2}\left\{\frac{e^{(-a-s)(t)}}{-a-s}\right\}\Bigg|_0^\infty\tag{assume s$>$a for integral to be finite}\\
                            &=\frac{1}{2}\left(\frac{1}{s-a}+\frac{1}{s+a}\right)\\
                            &=\frac{1}{2}\left(\frac{2s}{s^2-a^2}\right)\\
                           I&=\frac{s}{s^2-a^2}
\end{align*}
The option that has the same laplace transform as the one in the question is Option(A). So, it is the correct answer.\bigskip\\
\hfill\begin{minipage}{\linewidth}
    \begin{enumerate}
        \item[Q.15] The median of the following set of numbers: 154, 130, 144, 137, 156, 146, 138,
149, 160, 138 is:
    \item[(A)] $145.0$
    \item[(B)] $145.2$
    \item[(C)] $151.0$
    \item[(D)] $146.0$
    %\item[SOL] NEED TO WRITE SOL!!!
    \end{enumerate}   
\end{minipage}\bigskip\\
\underline{Solution}\medskip\\
To find the median, we need to arrange the data in ascending or descending order.\\
After arranging them in ascending order, the data is as follows:\medskip\\
\{130,137,138,138,144,146,149,154,156,160\}\medskip\\
The number of elements in the data set are 10.\medskip\\
Since it is an even number, the median (M) is defined as the average of the middle 2 terms (in this case, average of 5\textsuperscript{th} and 6\textsuperscript{th} terms).\medskip\\
Therefore, the median(M) is $\dfrac{144+146}{2} = 145.0$.\medskip\\
\large\underline{\textbf{C Implementation }}\medskip\\
\begin{minted}[frame=single,framesep=10pt]{c}
#include <stdio.h>

void swap(int*a,int i,int j);
int partition(int*a, int low, int high);
int quicksort(int*a,int low,int high);


int main(){
  int a[]={130,137,138,138,144,146,149,154,156,160};
  int q=sizeof(a)/sizeof(a[0]);
  //Printing original array
  printf("The original array is : [");
  for(int k=0;k<q;k++){
    printf("%d ",a[k]);
  }
  printf("]\n");
  //printf("%d\n", q);
  quicksort(a,0,q-1);
  printf("The sorted array is: [");
  for(int k=0;k<q;k++){
    printf("%d ",a[k]);
  }
  printf("]\n");
  double m;
  if (q%2==0){
    m=(a[(q/2)]+a[(q/2-1)])/2;
    printf("\nThe median is : %lf\n", m);
  }
  else {
  m=(a[(q-1)/2]);
  printf("\nThe median is : %lf\n", m);
  }
}

int partition(int* a,int low, int high){
  int pivot = a[high];
  int i=low-1;
  for (int j=low;j<high;j++){
    if(a[j]<pivot){
      i+=1;
      swap(a,i,j);
    }
  }
  swap(a,i+1,high);
  return i+1;
}
int quicksort(int *a, int low ,int high){
  if(low<high){
    int pi=partition(a,low,high);
  quicksort(a,low,pi-1);
  quicksort(a,pi+1,high);
  }
}

void swap(int*a,int i,int j){
  int temp=a[i];
  a[i]=a[j];
  a[j]=temp;
}
    
\end{minted}

\noindent The output is as follows:\medskip\\
\begin{lstlisting}[style=terminal]
madhur@madhur-Dell-15-DC15250:~/LATEX_CODES$ gcc Sorter.c
madhur@madhur-Dell-15-DC15250:~/LATEX_CODES$ ./a.out
The original array is : [130 137 138 138 144 146 149 154 156 160 ]
The sorted array is: [130 137 138 138 144 146 149 154 156 160 ]

The median is : 145.000000

\end{lstlisting}
\subsection*{Q.36 – Q.65 Carry TWO marks each}

\hfill\begin{minipage}{\linewidth}
    \begin{enumerate}
    \item[Q.36] The series \\ $\sum_{n=0}^{r} q^n$ $=1+q+q^2+... ... ...$ has the sum:
    \item[(A)] $\frac{1}{1-q}$-$\frac{q^{r+1}}{1-q}$
    \item[(B)] $\frac{1}{q-1}$+$\frac{q^{r+1}}{q-1}$
    \item[(C)] $\frac{1}{1-q}$+$\frac{q^{r+1}}{1-q}$
    \item[(D)] $\frac{1}{q-1}$
    %\item[SOL] NEED TO WRITE SOL!!! 
    \end{enumerate}
\end{minipage}
\underline{Solution}\medskip\\
    The given sequence is a geometric progression with initial term as 1 and common ratio as q and the number of terms are r+1.\\
    The sum of a geometrical progression with initial term a, common ratio r and number of terms n is given by:\\
    \begin{center}
        \large$a\,\frac{r^n-1}{r-1}$
    \end{center}
    applying the above formula, we get :
    \begin{center}
        $1\,\frac{q^{r+1}-1}{q-1}$\medskip\\
        $\frac{1-q^{r+1}}{1-q}$\medskip\\
        $\frac{1}{1-q}-\frac{q^{r+1}}{1-q}$
        
    \end{center}
Therefore, the answer is option A.\medskip\\
\large\underline{\textbf{C implementation}}\\

\begin{minted}[frame=single, framesep=10pt]{c}
#include<stdio.h>
#include<math.h>
int main(){
  double q,sum=0;
  int r;
  printf("Enter the value of q:");
  scanf("%lf",&q);
  printf("\nEnter the value of r:");
  scanf("%d",&r);
  for(int i=0;i<=r;i++){
    sum+=pow(q,i);
    
  }
  double formula =(1/(1-q))-((pow(q,r+1))/(1-q));
  printf("Sum is %lf\n", sum);
  printf("The formula gave %lf\n",formula);
  if (sum==formula){
  printf("Both are equal\n");
  }
}   
\end{minted}

\begin{lstlisting}[style=terminal]
madhur@madhur-Dell-15-DC15250:~/LATEX_CODES$ gcc gp.c -lm
madhur@madhur-Dell-15-DC15250:~/LATEX_CODES$ ./a.out
Enter the value of q:2

Enter the value of r:3
Sum is 15.000000
The formula gave 15.000000
Both are equal
\end{lstlisting}
    
\hfill\begin{minipage}{\linewidth}
    \begin{enumerate}
    \item[Q.37] The value of $\lim_{x\to 0}\dfrac{x\cos(x)-\sin(x)}{x^2\sin(x)}$
        \item[(A)] $\frac{-1}{3}$
        \item[(B)] $\frac{1}{3}$
        \item[(C)] $3$
        \item[(D)]$-3$
        %\item[SOL] NEED TO WRITE SOL!!! 
    \end{enumerate}
        \underline{Solution}\medskip\\
        We use the idea of taylor series expansion to solve the problem.\\
        
        \begin{center}
        $\frac{x(\cos(x))-\sin(x)}{x^2\sin(x)}$\\
        $\frac{x(1-\frac{x^2}{2!} .....)-(x-\frac{x^3}{3!}.....)}{x^2(x-\frac{x^3}{3!}.....)}$\\
        When x approaches zero, we can ignore the higher power terms.\\
        $\frac{(x-\frac{x^3}{2})-(x-\frac{x^3}{6})}{x^3(1-\frac{x^2}{6})}$\\
        $\frac{\frac{-1}{3}}{1-\frac{x^2}{6}}$\\
        Substituting x as zero gives us the limit as $\dfrac{-1}{3}$.
        \end{center}
        
\end{minipage}
\large\underline{\textbf{Python Implementation}}\bigskip\\
\begin{lstlisting}[style=pycode,frame=single]
import matplotlib.pyplot as plt
import numpy as np

def f(x):
    return (x*(np.cos(x))-np.sin(x))/((x**2)*np.sin(x))

x = np.linspace(-0.5, 0.5, 1000)
y = f(x)
plt.plot(x, y, color='blue', linestyle='-', linewidth=2, label='f(x) = Given Function')
plt.axhline(y=-1/3, color='red', linestyle='-', linewidth=2, label='f(x) = -0.33333')
plt.title('Func limit is -0.33333')
plt.xlabel('X Axis')
plt.ylabel('Y Axis')
plt.grid(True, alpha=0.5)
plt.legend() 
plt.show()
\end{lstlisting}\vspace{1cm}

  \large The output graph is shown below:\\
\begin{figure}[htbp]
    \centering
    \includegraphics[width=0.7\textwidth]{Images/limitplot.png}
\end{figure} \newpage 


\hfill\begin{minipage}{\linewidth}
\begin{enumerate}
    \item[Q.38] If $u=x^4+y^4+3x^2y^2$, the value of $\left(x\dfrac{\partial{u}}{\partial{x}}+y\dfrac{\partial{u}}{\partial{y}}\right)^{-1}$
    \item[(A)] $4u$
    \item[(B)] $0$
    \item[(C)] $\dfrac{1}{4u}$
    \item[(D)] $2u$
    %\item[SOL] NEED TO WRITE SOL!!! 
\end{enumerate}
\end{minipage}\bigskip


\underline{Solution}\\
We can directly evaluate the partial derivative using the chain and product rules, but another approach is to use euler's theorem for homogeneous functions.\\
We will solve it using the second method, since the first one is straightforward.\bigskip\\
Step 1: Proving the given function is homogeneous:\medskip\\
A function $f(x,y)$ is said to be homogeneous and having degree n if $f(\alpha x, \alpha y)=\alpha^n f(x,y)$\medskip\\
now, let 
\begin{align*}
    u=g(x,y)&=x^4+y^4+3x^2y^2\\
g(\alpha x,\alpha y)&=\alpha^4(x^4+y^4+3x^2y^2)\\    
\end{align*}
Therefore, u is a homogeneous function of degree 4.\medskip\\
Step 2: Using Euler's theorem to get the answer\medskip\\
\textbf{Euler's Theorem for homogeneous functions}\\
\quad It states that if a function $f(x,y)$ is homogeneous, then the following condition is satisfied:\begin{align*}
    x\frac{\partial u}{\partial x}+y\frac{\partial u}{\partial y}&=nf\\  
\end{align*}
Since we know that the given equation is homogeneous and of degree 4, we can apply this theorem.\\
\begin{align*}\
    x\frac{\partial u}{\partial x}+y\frac{\partial u}{\partial y}&=4u\\
\end{align*}
Since they asked for its reciprocal, the answer is $\dfrac{1}{4u}$\\

\hfill\begin{minipage}{\linewidth}
\begin{enumerate}
    \item[Q.43] The value of $y$ at $x=0.3$ and $y(x=0)=0$ for the differential equation $\frac{dy}{dx}=3y+2e^x$ is \underline{\hspace{1cm}}. (Rounded off to 3 decimal places)  
    %\item[SOL] NEED TO WRITE SOL!!! 
\end{enumerate}    
\end{minipage}\bigskip\\
\underline{Solution}\\\\
Step 1: Defining a function G(s) using laplace transform:\\ 
\begin{align*}
    G(s)&=\mathcal{L}\{y\}\\
    sG(s)&=\mathcal{L}\{y'\} \tag{By property of laplace transform and Given y(0)=0}
\end{align*}
 Step 2: Applying the laplace transform on the diiferential equation:\\
 \begin{align*}
     \mathcal{L}\{y'\}&=\mathcal{L}\{3y+2e^x\}\\
     \mathcal{L}\{y'\}&=3\mathcal{L}\{y\}\,+2\mathcal{L}\{e^x\}\tag{Using linearity of laplace transform}\\
     sG(s)&=3G(s)+2\mathcal{L}\{e^x\}
 \end{align*}
Step 3:Finding laplace transform of $e^x$:\\
\begin{align*}
    \mathcal{L}\{e^x\}&=\int_{0}^{\infty}e^t\,e^{-st} dt\\
                      &=\int_{0}^{\infty}e^{(1-s)(t)} dt\\
                      &=\frac{1}{1-s}\,\,(e^{(1-s)(t)}\big|_{0}^{\infty})\\
                      &=\frac{1}{1-s}(0-1)\tag{Assuming that s$>$1}\\
                      &=\frac{1}{s-1}
\end{align*}
Step 4:Finding G(s):\\
\begin{align*}
    G(s)(s-3)&=\frac{2}{s-1}\\
    G(s)&=\frac{2}{(s-1)(s-3)}\\
    \text{Using partial fractions:}\\
    G(s)&=\left(\frac{1}{s-3}-\frac{1}{s-1}\right)
\end{align*}
Step 5: Using inverse laplace transform to find y:
\begin{align*}
   \mathcal{L}\{y\}&=G(s)\\
                  y &=\mathcal{L^{-1}}\left\{\frac{1}{s-3}-\frac{1}{s-1}\right\}\\
                  y &=\mathcal{L^{-1}}\left\{\frac{1}{s-3}\right\}-\mathcal{L^{-1}}\left\{\frac{1}{s-1}\right\}\tag{Using linearity of inverse laplace transform}\\
                  y&=e^{3x}-e^{x}\tag{From above, we say that$\mathcal{L}\{e^{ax}\}=\frac{1}{s-a}$ provided s$>$a.}\\
\end{align*}
Step 6: Substituting condition to find value of y:\\
\begin{align*}
    x&=0.3\tag{given}\\
    y&=e^{3\times0.3}-e^{0.3}\\
    y&=1.110
\end{align*}
\large\underline{\textbf{Python Implementation}}\\
\begin{lstlisting}[style=pycode,frame=single]
# The expression is y=e^{3x}-e^{x}
#evaluating the expression:
import math
x=0.3
y=math.exp(3*x)-math.exp(x)
print(y)
\end{lstlisting}
Output :\\
\begin{lstlisting}[style=terminal]
1.1097443035809462
\end{lstlisting}



\end{document}
